\documentclass[../../main.tex]{subfiles}% 注意这里的文件路径不能用 ./main.tex ,否则用latexmk编译子文件会报错
\graphicspath{{\subfix{./image/}}} % 指定图片目录,后续可以直接使用图片文件名
% 注意这里的文件路径不能用 ../../image/ ,否则用latexmk编译子文件会报错

% 例如:
% \begin{figure}[H]
% \centering
% \includegraphics[scale=0.3]{图.png}
% \caption{}
% \label{figure:图}
% \end{figure}
% 注意:上述\label{}一定要放在\caption{}之后,否则引用图片序号会只会显示??.

\begin{document}

\section{Gauss定理}

\begin{definition}
设正则曲面$S$的第一、第二基本形式系数为$g_{\alpha\beta},b_{\alpha\beta}$，定义\textbf{Gauss曲率}
\begin{align*}
K = \frac{b_{11}b_{22}-(b_{12})^2}{g_{11}g_{22}-(g_{12})^2} = \kappa_1\kappa_2,
\end{align*}
其中$\kappa_1,\kappa_2$为曲面的主曲率。
\end{definition}

\begin{proposition}
正则曲面的Gauss曲率满足
\begin{align}
K = \frac{R_{1212}}{g_{11}g_{22}-(g_{12})^2}. \label{eq:gauss_curv_riemann}
\end{align}
\end{proposition}

\begin{proof}
由Gauss方程$b_{11}b_{22}-(b_{12})^2=R_{1212}$，等式两边同时除以$g_{11}g_{22}-(g_{12})^2$，即得式\eqref{eq:gauss_curv_riemann}。

\end{proof}

\begin{theorem}
曲面的Gauss曲率是曲面在\textbf{保长变换}下的不变量。
\end{theorem}

\begin{proof}
Riemann记号$R_{1212}$仅由曲面第一基本形式及其一、二阶偏导数决定，结合式\eqref{eq:gauss_curv_riemann}可知$K$仅依赖第一基本形式。
保长变换保持曲面第一基本形式不变，故Gauss曲率在保长变换下不变。

\end{proof}

\begin{remark}
Gauss绝妙定理开创了微分几何的新纪元。仅由第一基本形式决定的曲面几何性质称为\textbf{内蕴几何}。将内蕴几何推广至$n$维空间，即在区域上给定正定对称二次微分形式$\mathrm{d}s^2=g_{\alpha\beta}\mathrm{d}u^\alpha\mathrm{d}u^\beta$并研究其弯曲性质，即为\textbf{Riemann几何}。
球面的Gauss曲率为正常数，平面的Gauss曲率为$0$，因此球面与平面之间不存在保长对应。
\end{remark}

\begin{proposition}
\begin{enumerate}[(1)]
\item 若曲面取正交参数曲线网$(u,v)$（$F=0$），则
\begin{align}
K = -\frac{1}{\sqrt{EG}}\left(\left(\frac{(\sqrt{E})_v}{\sqrt{G}}\right)_v+\left(\frac{(\sqrt{G})_u}{\sqrt{E}}\right)_u\right). \label{eq:gauss_curv_ortho}
\end{align}
\item 若曲面取等温参数系$(u,v)$，即$\mathrm{I}=\lambda^2\left((\mathrm{d}u)^2+(\mathrm{d}v)^2\right)$，则
\begin{align}
K = -\frac{1}{\lambda^2}\left(\frac{\partial^2}{\partial u^2}+\frac{\partial^2}{\partial v^2}\right)\ln\lambda. \label{eq:gauss_curv_isotherm}
\end{align}
\end{enumerate}
\end{proposition}

\begin{proof}
将正交参数下$R_{1212}$的表达式代入式\eqref{eq:gauss_curv_riemann}化简，可得式\eqref{eq:gauss_curv_ortho}；
将等温参数的第一基本形式代入式\eqref{eq:gauss_curv_ortho}，化简可得式\eqref{eq:gauss_curv_isotherm}。

\end{proof}

\begin{theorem}
$\mathbb{E}^3$中无脐点的曲面$S$为可展曲面的充分必要条件是其Gauss曲率$K\equiv0$。
\end{theorem}
\begin{proof}
{\heiti 必要性}: 可展曲面局部可与平面区域建立保长对应，平面的Gauss曲率恒为$0$，由Gauss绝妙定理，可展曲面的Gauss曲率$K\equiv0$。

{\heiti 充分性}: 设曲面$S$的Gauss曲率$K\equiv0$，任取$p\in S$，存在$p$的邻域$U$，其上可取正交曲率线网$(u,v)$，满足$F=M=0$，且$K=\frac{LN}{EG}\equiv0$。不妨设$N=0,L\neq0$。
由Codazzi方程$\frac{\partial N}{\partial u}=H\frac{\partial G}{\partial u}$，其中平均曲率$H=\frac{L}{2E}\neq0$，故$\frac{\partial G}{\partial u}=0$。
由自然标架运动公式
\begin{align*}
\boldsymbol{r}_{vv}=\Gamma^1_{22}\boldsymbol{r}_u+\Gamma^2_{22}\boldsymbol{r}_v+N\boldsymbol{n}=\Gamma^1_{22}\boldsymbol{r}_u+\Gamma^2_{22}\boldsymbol{r}_v,
\end{align*}
正交参数下$\Gamma^1_{22}=-\frac{1}{2E}\frac{\partial G}{\partial u}\equiv0$，因此$\boldsymbol{r}_{vv}\times\boldsymbol{r}_v=\boldsymbol{0}$，即$v$-曲线为直线，曲面为直纹面。
又$\boldsymbol{n}_v\cdot\boldsymbol{r}_u=-M=0,\boldsymbol{n}_v\cdot\boldsymbol{r}_v=-N=0,\boldsymbol{n}_v\cdot\boldsymbol{n}=0$，故$\boldsymbol{n}_v=\boldsymbol{0}$，即单位法向量沿$v$-曲线不变，因此$S$为可展曲面。

\end{proof}

\begin{theorem}
无脐点的曲面$S$为可展曲面的充分必要条件是$S$可与平面区域建立保长对应。
\end{theorem}
\begin{proof}
{\heiti 必要性}: 第三章已证可展曲面局部可与平面区域建立保长对应。

{\heiti 充分性}: 若$S$可与平面区域建立保长对应，由Gauss绝妙定理，$S$的Gauss曲率恒为$0$，结合上一定理，$S$为可展曲面。

\end{proof}

\begin{example}
设曲面$S,\tilde{S}$的参数方程为
\begin{align*}
\boldsymbol{r} &= \left(au,bv,\frac{1}{2}(au^2+bv^2)\right),\\
\tilde{\boldsymbol{r}} &= \left(\tilde{a}\tilde{u},\tilde{b}\tilde{v},\frac{1}{2}(\tilde{a}\tilde{u}^2+\tilde{b}\tilde{v}^2)\right),
\end{align*}
且$ab=\tilde{a}\tilde{b}$。证明：在对应$\tilde{u}=u,\tilde{v}=v$下，$S$与$\tilde{S}$有相同Gauss曲率；当$(a^2,b^2)\neq(\tilde{a}^2,\tilde{b}^2)$且$(a^2,b^2)\neq(b^2,a^2)$时，二者无保长对应。
\end{example}

\begin{proof}
直接计算曲面$S$的第一、第二基本形式：
\begin{align*}
\mathrm{I} &= a^2(1+u^2)(\mathrm{d}u)^2+2abuv\mathrm{d}u\mathrm{d}v+b^2(1+v^2)(\mathrm{d}v)^2,\\
\mathrm{II} &= \frac{a}{\sqrt{1+u^2+v^2}}(\mathrm{d}u)^2+\frac{b}{\sqrt{1+u^2+v^2}}(\mathrm{d}v)^2.
\end{align*}
代入Gauss曲率定义得
\begin{align}
K = \frac{1}{ab(1+u^2+v^2)^2}. \label{eq:ex_curv_S}
\end{align}
同理可得曲面$\tilde{S}$的Gauss曲率
\begin{align*}
\tilde{K} = \frac{1}{\tilde{a}\tilde{b}(1+\tilde{u}^2+\tilde{v}^2)^2} = \frac{1}{ab(1+u^2+v^2)^2}=K,
\end{align*}
故对应点处二者Gauss曲率相同。

假设存在保长对应$\tilde{u}=\tilde{u}(u,v),\tilde{v}=\tilde{v}(u,v)$，由Gauss绝妙定理，对应点Gauss曲率相等，即
\begin{align*}
\frac{1}{ab(1+\tilde{u}^2+\tilde{v}^2)^2}=\frac{1}{ab(1+u^2+v^2)^2},
\end{align*}
化简得$\tilde{u}^2+\tilde{v}^2=u^2+v^2$。
在$(u,v)=(0,0)$处，有$\tilde{u}(0,0)=\tilde{v}(0,0)=0$。对$\tilde{u}^2+\tilde{v}^2=u^2+v^2$求一阶偏导得
\begin{align*}
\begin{cases}
\tilde{u}\frac{\partial\tilde{u}}{\partial u}+\tilde{v}\frac{\partial\tilde{v}}{\partial u}=u,\\
\tilde{u}\frac{\partial\tilde{u}}{\partial v}+\tilde{v}\frac{\partial\tilde{v}}{\partial v}=v.
\end{cases}
\end{align*}
再次求二阶偏导并代入$(u,v)=(0,0)$，得Jacobi矩阵
\begin{align*}
J=\begin{pmatrix}
\frac{\partial\tilde{u}}{\partial u}&\frac{\partial\tilde{v}}{\partial u}\\
\frac{\partial\tilde{u}}{\partial v}&\frac{\partial\tilde{v}}{\partial v}
\end{pmatrix}
\end{align*}
在原点处为正交矩阵，设
\begin{align*}
J|_{(0,0)}=\begin{pmatrix}
\cos\theta&\sin\theta\\
-\varepsilon\sin\theta&\varepsilon\cos\theta
\end{pmatrix},\quad \varepsilon=\pm1.
\end{align*}
由保长对应性质$\mathrm{I}=J\tilde{\mathrm{I}}J^\mathrm{T}$，代入原点处得
\begin{align*}
\begin{pmatrix}a^2&0\\0&b^2\end{pmatrix}=
\begin{pmatrix}\cos\theta&\sin\theta\\-\varepsilon\sin\theta&\varepsilon\cos\theta\end{pmatrix}
\begin{pmatrix}\tilde{a}^2&0\\0&\tilde{b}^2\end{pmatrix}
\begin{pmatrix}\cos\theta&-\varepsilon\sin\theta\\\sin\theta&\varepsilon\cos\theta\end{pmatrix},
\end{align*}
展开得方程组
\begin{align*}
\begin{cases}
\tilde{a}^2\cos^2\theta+\tilde{b}^2\sin^2\theta=a^2,\\
(\tilde{b}^2-\tilde{a}^2)\sin\theta\cos\theta=0,\\
\tilde{a}^2\sin^2\theta+\tilde{b}^2\cos^2\theta=b^2.
\end{cases}
\end{align*}
若$\tilde{b}^2\neq\tilde{a}^2$，则$\theta=0$或$\theta=\frac{\pi}{2}$，即$(a^2,b^2)=(\tilde{a}^2,\tilde{b}^2)$或$(a^2,b^2)=(\tilde{b}^2,\tilde{a}^2)$。
因此当$(a^2,b^2)\neq(\tilde{a}^2,\tilde{b}^2)$且$(a^2,b^2)\neq(\tilde{b}^2,\tilde{a}^2)$时，曲面$S$与$\tilde{S}$之间不存在保长对应。

\end{proof}

\begin{theorem}
设$\sigma:S_1\to S_2$为连续可微映射，$S_1$无脐点且Gauss曲率$K\neq0$。若$S_1,S_2$在所有对应点、对应切方向的法曲率相等，则存在$\mathbb{E}^3$上的刚体运动$\tilde{\sigma}$，使得$\sigma=\tilde{\sigma}|_{S_1}$。
\end{theorem}
\begin{proof}
因$S_1$无脐点，可取正交曲率线网$(u,v)$，则$S_1$的基本形式为
\begin{align*}
\mathrm{I}=E(\mathrm{d}u)^2+G(\mathrm{d}v)^2,\quad \mathrm{II}=L(\mathrm{d}u)^2+N(\mathrm{d}v)^2,
\end{align*}
且$L=\kappa_1E,N=\kappa_2G,\kappa_1>\kappa_2,K=\kappa_1\kappa_2\neq0$。
由法曲率对应相等的条件，$(u,v)$为$S_2$的正交曲率线网，故$S_2$的基本形式满足
\begin{align*}
\tilde{F}=\tilde{M}=0,\quad \tilde{L}=\kappa_1\tilde{E},\quad \tilde{N}=\kappa_2\tilde{G}.
\end{align*}
法曲率相等条件为
\begin{align*}
\frac{\kappa_1E(\mathrm{d}u)^2+\kappa_2G(\mathrm{d}v)^2}{E(\mathrm{d}u)^2+G(\mathrm{d}v)^2}=
\frac{\kappa_1\tilde{E}(\mathrm{d}u)^2+\kappa_2\tilde{G}(\mathrm{d}v)^2}{\tilde{E}(\mathrm{d}u)^2+\tilde{G}(\mathrm{d}v)^2},
\end{align*}
展开化简得$(\kappa_1-\kappa_2)(\tilde{E}G-E\tilde{G})=0$，由$\kappa_1-\kappa_2>0$，得$\tilde{E}G=E\tilde{G}$。
设$\frac{\tilde{E}}{E}=\frac{\tilde{G}}{G}=\lambda$，则$\tilde{\mathrm{I}}=\lambda\mathrm{I},\tilde{\mathrm{II}}=\lambda\mathrm{II}$。

由正交曲率线网下的Codazzi方程，$S_1$满足
\begin{align*}
\frac{\partial L}{\partial v}=H\frac{\partial E}{\partial v},\quad \frac{\partial N}{\partial u}=H\frac{\partial G}{\partial u},\quad H=\frac{1}{2}(\kappa_1+\kappa_2),
\end{align*}
$S_2$满足
\begin{align*}
\frac{\partial\tilde{L}}{\partial v}=H\frac{\partial\tilde{E}}{\partial v},\quad \frac{\partial\tilde{N}}{\partial u}=H\frac{\partial\tilde{G}}{\partial u}.
\end{align*}
代入$\tilde{E}=\lambda E,\tilde{G}=\lambda G,\tilde{L}=\lambda L,\tilde{N}=\lambda N$，化简得$(\kappa_1-\kappa_2)\frac{\partial\lambda}{\partial u}=(\kappa_1-\kappa_2)\frac{\partial\lambda}{\partial v}=0$，故$\lambda$为常数。

由正交参数下Gauss曲率公式，$\tilde{K}=\frac{1}{\lambda}K$，结合$\tilde{K}=K\neq0$，得$\lambda=1$。
因此$S_1,S_2$的第一、第二基本形式完全相同，由曲面基本形式的存在唯一性定理，存在$\mathbb{E}^3$上的刚体运动$\tilde{\sigma}$，使得$\sigma=\tilde{\sigma}|_{S_1}$。

\end{proof}

\begin{example}
已知曲面$S$的第一基本形式如下，求它们的Gauss曲率$K$。
\begin{enumerate}[(1)]
\item $\mathrm{I}=\frac{(\mathrm{d}u)^2+(\mathrm{d}v)^2}{\left(1+\frac{c}{4}(u^2+v^2)\right)^2}$，$c$是常数；
\item $\mathrm{I}=\frac{a^2\left((\mathrm{d}u)^2+(\mathrm{d}v)^2\right)}{v^2}$，$v>0$，$a$是常数；
\item $\mathrm{I}=\frac{(\mathrm{d}u)^2+(\mathrm{d}v)^2}{(u^2+v^2+c^2)^2}$，$c>0$是常数；
\item $\mathrm{I}=(\mathrm{d}u)^2+\mathrm{e}^{\frac{2u}{a}}(\mathrm{d}v)^2$，$a$是常数；
\item $\mathrm{I}=(\mathrm{d}u)^2+\cosh^2\frac{u}{a}(\mathrm{d}v)^2$，$a$是常数；
\item $\mathrm{I}=(\mathrm{d}u)^2+G(u)(\mathrm{d}v)^2$.
\end{enumerate}
\end{example}

\begin{example}
证明在下列曲面之间不存在保长对应:
\begin{enumerate}[(1)]
\item 球面；
\item 柱面；
\item 双曲抛物面: $z=x^2-y^2$.
\end{enumerate}
\end{example}

\begin{example}
已知曲面$S$和$\tilde{S}$的第一基本形式分别是
\begin{align*}
\mathrm{I} &= (\mathrm{d}u)^2+(1+u^2)(\mathrm{d}v)^2,\\
\tilde{\mathrm{I}} &= \frac{\tilde{u}^2}{\tilde{u}^2-1}\mathrm{d}\tilde{u}^2+\tilde{u}^2\mathrm{d}\tilde{v}^2,
\end{align*}
试问: 在曲面$S$和$\tilde{S}$之间是否存在保长对应?
\end{example}

\begin{example}
设曲面$S$和$\tilde{S}$的参数方程分别是
\begin{align*}
\boldsymbol{r} &= (u\cos v,u\sin v,\ln u),\\
\tilde{\boldsymbol{r}} &= (\tilde{u}\cos\tilde{v},\tilde{u}\sin\tilde{v},\tilde{v}).
\end{align*}
证明: 在对应$u=\tilde{u}$, $v=\tilde{v}$下, 这两个曲面在对应点有相同的Gauss曲率, 但是在曲面$S$和$\tilde{S}$之间不存在保长对应.
\end{example}

\begin{example}
已知曲面$S$和$\tilde{S}$的第一基本形式分别是
\begin{align*}
\mathrm{I} &= \mathrm{e}^{2v}\left((\mathrm{d}u)^2+a^2(1+u^2)(\mathrm{d}v)^2\right),\\
\tilde{\mathrm{I}} &= \mathrm{e}^{2\tilde{v}}\left((\mathrm{d}\tilde{u})^2+b^2(1+\tilde{u}^2)(\mathrm{d}\tilde{v})^2\right),
\end{align*}
其中$a^2\neq b^2$. 证明: 在对应$u=\tilde{u}$, $v=\tilde{v}$下, 这两个曲面在对应点有相同的Gauss曲率, 但是该对应不是曲面$S$和$\tilde{S}$之间的保长对应.
\end{example}

\begin{example}
设曲面$S$的第一基本形式是$\mathrm{I}=E(\mathrm{d}u)^2+2F\mathrm{d}u\mathrm{d}v+G(\mathrm{d}v)^2$. 证明: 它的Gauss曲率的表达式是
\begin{align*}
K &= \frac{1}{(EG-F^2)^2}
\left(
\begin{vmatrix}
-\frac{G_{uu}}{2}+F_{uv}-\frac{E_{vv}}{2} & \frac{E_u}{2} & F_u-\frac{E_v}{2}\\[4pt]
F_v-\frac{G_u}{2} & E & F\\[4pt]
\frac{G_v}{2} & F & G
\end{vmatrix}
-
\begin{vmatrix}
0 & \frac{E_v}{2} & \frac{G_u}{2}\\[4pt]
\frac{E_v}{2} & E & F\\[4pt]
\frac{G_u}{2} & F & G
\end{vmatrix}
\right).
\end{align*}
\end{example}

\begin{example}
若在定理5.4中曲面$S_1$的Gauss曲率$K=0$, 则定理的结论是否成立? 举例说明.
\end{example}




\end{document}